\section{Block-diagonal boundary conditions for periodic ground spaces}
\label{sec:block_diagonal_boundary_closing}

For a block-diagonal tensor $B=\bigoplus_{j=1}^g\mu_jA_j$ with $\mu_j\ne0$,
simultaneous injectivity at length $S>0$ gives the periodic ground-space
identity at every $S+1\leq L\leq N$, including the minimal ring
$N=S+1$; see Theorem~\ref{thm:ph_block_ground_space_at_simultaneous_injectivity}.
The earlier separation argument starts from individual block injectivity
and uses a larger interaction range. Its explicit boundary identities yield
component membership
$\Gamma_N^{A_j}(\mu_j^NX_j)\in\mathcal G_{N,L}(A_j)$, establishing the
periodic boundary step at that larger range. The simultaneous-span argument
gives the bound of \cite[Theorem~IV.16]{Cirac2021Matrix} directly.

\begin{lemma}[Normalized BNT representatives give the large-length block intersection]
    \label{lem:bnt_direct_sum_unital_block_intersection_ge}
    \lean{MPSTensor.pgvwc07_directSum_restriction_intersection_of_wordTupleSpanTop}
    \lean{MPSTensor.pgvwc07_iSup_restriction_intersection_of_ge_of_bnt_directSum_unital_c1}
    \lean{MPSTensor.pgvwc07_directSum_restriction_intersection_of_ge_of_bnt_directSum_unital_c1}
    \leanok
    \uses{lem:unital_block_separating_product_span, lem:product_span_block_image_direct_sum, lem:block_restriction_intersection_subspace}
    Let $A^1,\ldots,A^r$ be normalized representatives of a basis of normal
    tensors: each representative is irreducible, the family is left-canonical,
    the self-overlaps are normalized, and no two distinct same-dimensional
    representatives are gauge-phase equivalent. Suppose $L_0>0$, and, for
    every block $j$, the map
    \begin{align}
        \Gamma_{L_0}^{A^j}:\MN{D_j}
         & \longrightarrow(\C^d)^{\otimes L_0}
        \notag
    \end{align}
    is injective. Assume also the normalization
    $\sum_a A^j_aA^{j\,\dagger}_a=\Id$. Set
    $N=(L_0+1)+(r-1)3(L_0+1)$. For every $n\ge N$, set
    $S_n:=\bigvee_jG_{n+1}(A^j)$. The two consecutive sums are internal direct
    sums:
    \begin{align}
        \bigvee_jG_{n+1}(A^j)
         & =\bigoplus_jG_{n+1}(A^j),
        \notag                       \\
        \bigvee_jG_{n+2}(A^j)
         & =\bigoplus_jG_{n+2}(A^j).
        \notag
    \end{align}
    With this notation,
    \begin{align}
         & \left\{\psi:\forall b,\ \psi(-,b)\in S_n\right\}
        \cap
        \left\{\psi:\forall a,\ \psi(a,-)\in S_n\right\} =\bigvee_jG_{n+2}(A^j).
        \notag
    \end{align}
\end{lemma}

\begin{proof}
    \leanok
    Lemma~\ref{lem:unital_block_separating_product_span} gives
    $\spn\{(A^1_w,\ldots,A^r_w):|w|=n\} =\prod_j\MN{D_j}$
    for every $n\ge N$. Substituting this product-span equation into
    Lemma~\ref{lem:block_restriction_intersection_subspace}, together with the
    normalization equation, gives the displayed intersection identity.
    Lemma~\ref{lem:product_span_block_image_direct_sum} at lengths $n+1$ and
    $n+2$ gives the two directness statements.
\end{proof}

\begin{lemma}[Normalized BNT representatives give an eventual block intersection]
    \label{lem:bnt_direct_sum_unital_block_intersection}
    \lean{MPSTensor.pgvwc07_iSup_restriction_intersection_eventually_of_bnt_directSum_unital_c1}
    \leanok
    \uses{lem:bnt_direct_sum_unital_block_intersection_ge}
    Let $A^1,\ldots,A^r$ be normalized representatives of a basis of normal
    tensors: each representative is irreducible, the family is left-canonical,
    the self-overlaps are normalized, and no two distinct same-dimensional
    representatives are gauge-phase equivalent. Suppose $L_0>0$, and, for
    every block $j$, the map
    \begin{align}
        \Gamma_{L_0}^{A^j}:\MN{D_j}
         & \longrightarrow(\C^d)^{\otimes L_0}
        \notag
    \end{align}
    is injective. Assume also the normalization
    $\sum_a A^j_aA^{j\,\dagger}_a=\Id$. Then there is a length $N$ such that,
    for every $n\ge N$, with $S_n:=\bigvee_jG_{n+1}(A^j)$, one has
    \begin{align}
         & \left\{\psi:\forall b,\ \psi(-,b)\in S_n\right\}
        \cap
        \left\{\psi:\forall a,\ \psi(a,-)\in S_n\right\} =\bigvee_jG_{n+2}(A^j).
        \notag
    \end{align}
\end{lemma}

\begin{proof}
    \leanok
    Apply Lemma~\ref{lem:bnt_direct_sum_unital_block_intersection_ge} and take
    $N=(L_0+1)+(r-1)3(L_0+1)$.
\end{proof}

\begin{lemma}[Block-diagonal boundary-condition formula]
    \label{lem:block_diagonal_boundary_condition_sum}
    \lean{MPSTensor.BlockSumGroundSpace.groundSpaceMap_toTensorFromBlocks_eq_sum_blockDiagonal}
    \leanok
    \uses{def:ground_space_parent}
    Let $B=\bigoplus_{j=1}^g\mu_jA_j$. For block-diagonal boundary conditions
    $X_j$,
    \begin{align}
        \Gamma_L^B\!\left(\bigoplus_{j=1}^g X_j\right)
         & =\sum_{j=1}^g\Gamma_L^{A_j}(\mu_j^LX_j).
        \notag
    \end{align}
\end{lemma}

\begin{proof}
    \leanok
    This is the diagonal-block trace formula specialized to a block-diagonal
    boundary matrix. The $j$-th diagonal block of $\bigoplus_kX_k$ is $X_j$,
    giving the displayed sum.
\end{proof}

\begin{lemma}[Boundary matrices for a sum of open-boundary block spaces]
    \label{lem:block_diagonal_boundary_of_open_boundary_sum}
    \lean{MPSTensor.BlockSumGroundSpace.exists_blockDiagonal_boundary_of_mem_iSup_groundSpace}
    \leanok
    \uses{def:ground_space_parent, lem:block_diagonal_boundary_condition_sum}
    Let $B=\bigoplus_{j=1}^g\mu_jA_j$, where $\mu_j\ne0$ for every $j$.
    For every length $L$ and every
    $\psi\in\bigvee_{j=1}^gG_L(A_j)$, there are boundary matrices $X_j$ such
    that
    \begin{align}
        \psi
         & =\Gamma_L^B\!\left(\bigoplus_{j=1}^gX_j\right).
        \notag
    \end{align}
    For every $j$, one also has
    $\Gamma_L^{A_j}(\mu_j^LX_j)\in G_L(A_j)$.
\end{lemma}

\begin{proof}
    \leanok
    \uses{def:ground_space_parent, lem:block_diagonal_boundary_condition_sum}
    Choose $\phi_j\in G_L(A_j)$ such that
    \begin{align}
        \psi
         & =\sum_{j=1}^g\phi_j.
        \notag
    \end{align}
    Choose $Y_j$ with $\phi_j=\Gamma_L^{A_j}(Y_j)$, and set
    $X_j=\mu_j^{-L}Y_j$. Since
    $\mu_j\ne0$, one has, for every $j$,
    \begin{align}
        \Gamma_L^{A_j}(\mu_j^LX_j)
         & =\phi_j.
        \notag
    \end{align}
    Lemma~\ref{lem:block_diagonal_boundary_condition_sum} now gives
    \begin{align}
        \Gamma_L^B\!\left(\bigoplus_{j=1}^gX_j\right)
         & =\sum_{j=1}^g\Gamma_L^{A_j}(\mu_j^LX_j)
        =\sum_{j=1}^g\phi_j
        =\psi.
        \notag
    \end{align}
\end{proof}

\begin{lemma}[One block lies in the local parent space of the direct-sum tensor]
    \label{lem:local_ground_space_block_le_assembled}
    \lean{MPSTensor.groundSpace_block_le_toTensorFromBlocks}
    \leanok
    \uses{def:ground_space_parent}
    Let $B=\bigoplus_{k=1}^g\mu_kA_k$. If $\mu_j\ne0$, then, for every length
    $L$, one has $G_L(A_j)\subseteq G_L(B)$.
\end{lemma}

\begin{proof}
    \leanok
    A vector in $G_L(A_j)$ has the form $\Gamma_L^{A_j}(X)$. Put
    $\mu_j^{-L}X$ in the $j$-th diagonal block of a boundary matrix for $B$ and
    put zero in every other block. If $\iota_j(Z)$ denotes this block
    insertion, Lemma~\ref{lem:block_diagonal_boundary_condition_sum} gives
    \begin{align}
        \Gamma_L^B\!\left(\iota_j(\mu_j^{-L}X)\right)
         & =
        \sum_{k=1}^g
        \Gamma_L^{A_k}
        \!\left(\mu_k^L(\iota_j(\mu_j^{-L}X))_{kk}\right) =\Gamma_L^{A_j}(X).
        \notag
    \end{align}
    Hence $\Gamma_L^{A_j}(X)\in G_L(B)$.
\end{proof}

\begin{theorem}[Local ground space of a block-diagonal tensor]
    \label{thm:block_sum_local_ground_space}
    \lean{MPSTensor.groundSpace_toTensorFromBlocks_eq_iSup,
        MPSTensor.groundSpaceES_toTensorFromBlocks_eq_iSup}
    \leanok
    \uses{def:ground_space_parent, lem:block_diagonal_boundary_condition_sum, lem:local_ground_space_block_le_assembled}
    Let $B=\bigoplus_{j=1}^g\mu_jA_j$, where $\mu_j\ne0$ for
    $1\le j\le g$. Then, for every length $L$,
    $G_L(B)=\bigvee_{j=1}^gG_L(A_j)$.
\end{theorem}

\begin{proof}
    \leanok
    A boundary matrix $X$ for $B$ has diagonal blocks $X_j$, and
    block-diagonality gives
    $\Gamma_L^B(X) =\sum_{j=1}^g\Gamma_L^{A_j}(\mu_j^LX_j)$.
    This proves $G_L(B)\subseteq\bigvee_jG_L(A_j)$. The reverse inclusion is
    Lemma~\ref{lem:local_ground_space_block_le_assembled} for each block $j$.
\end{proof}

\begin{lemma}[Periodic constraints for a block-diagonal tensor]
    \label{lem:block_diagonal_periodic_constraints_propagated_sum}
    \lean{MPSTensor.chainGroundSpace_toTensorFromBlocks_le_iSup_groundSpace}
    \leanok
    \uses{thm:block_sum_local_ground_space, lem:periodic_constraints_propagated_subspace_sequence}
    Let $B=\bigoplus_{j=1}^g\mu_jA_j$, where $\mu_j\ne0$ for
    $1\le j\le g$, and set $S_M:=\bigvee_{j=1}^gG_M(A_j)$. If, for every
    $M\ge L$,
    $\C^d\otimes S_M\cap S_M\otimes\C^d =S_{M+1}$,
    then $\mathcal G_{N,L}(B)\subseteq S_N$.
\end{lemma}

\begin{proof}
    \leanok
    Theorem~\ref{thm:block_sum_local_ground_space} gives $G_L(B)=S_L$.
    Apply Lemma~\ref{lem:periodic_constraints_propagated_subspace_sequence}
    to the sequence $S_M$.
\end{proof}

\begin{lemma}[Periodic constraints lie in the block local-space sum]
    \label{lem:bnt_block_diagonal_periodic_constraints_propagated_sum}
    \lean{MPSTensor.chainGroundSpace_toTensorFromBlocks_le_iSup_groundSpace_of_ge_of_bnt_directSum_unital_c1}
    \leanok
    \uses{lem:bnt_direct_sum_unital_block_intersection_ge, lem:block_diagonal_periodic_constraints_propagated_sum}
    Let $A_1,\ldots,A_g$ be normalized representatives of a basis of normal
    tensors: each representative is irreducible, the family is left-canonical,
    the self-overlaps are normalized, and no two distinct same-dimensional
    representatives are gauge-phase equivalent. Suppose $L_0>0$, and, for
    every block $j$, the map
    \begin{align}
        \Gamma_{L_0}^{A_j}:\MN{D_j}
         & \longrightarrow(\C^d)^{\otimes L_0}
        \notag
    \end{align}
    is injective. Assume also the normalization
    $\sum_a A^j_aA^{j\,\dagger}_a=\Id$. For nonzero scalars $\mu_j$, set
    \begin{align}
        B   & =\bigoplus_{j=1}^g\mu_jA_j,
        \notag                            \\
        S_M & :=\bigvee_{j=1}^gG_M(A_j).
        \notag
    \end{align}
    If $L\ge (L_0+1)+(g-1)3(L_0+1)+1$, then, for every $N\ge L$,
    $\mathcal G_{N,L}(B)\subseteq S_N$.
\end{lemma}

\begin{proof}
    \leanok
    Put $T=(L_0+1)+(g-1)3(L_0+1)$. For $M\ge L$, one has $M-1\ge T$.
    Applying Lemma~\ref{lem:bnt_direct_sum_unital_block_intersection_ge} at
    $n=M-1$ gives
    $\C^d\otimes S_M\cap S_M\otimes\C^d =S_{M+1}$
    for every $M\ge L$.
    Lemma~\ref{lem:block_diagonal_periodic_constraints_propagated_sum} then
    propagates the local constraint from $L$ to $N$.
\end{proof}

\begin{lemma}[Periodic constraints and internality of the block local-space sum]
    \label{lem:bnt_block_diagonal_periodic_constraints_internal_sum}
    \lean{MPSTensor.chainGroundSpace_toTensorFromBlocks_le_iSup_and_iSupIndep_of_bnt_unital_c1}
    \leanok
    \uses{lem:bnt_block_diagonal_periodic_constraints_propagated_sum, lem:product_span_block_image_direct_sum}
    With the hypotheses and notation of
    Lemma~\ref{lem:bnt_block_diagonal_periodic_constraints_propagated_sum}, set
    $S_N:=\bigvee_{j=1}^gG_N(A_j)$. Then, for every $N\ge L$,
    $\mathcal G_{N,L}(B)\subseteq S_N$, and the sum defining $S_N$ is internal:
    if
    \begin{align}
        \phi_j             & \in G_N(A_j),
        \notag                             \\
        \sum_{j=1}^g\phi_j & =0,
        \notag
    \end{align}
    then $\phi_j=0$ for every $j$.
\end{lemma}

\begin{proof}
    \leanok
    The inclusion is
    Lemma~\ref{lem:bnt_block_diagonal_periodic_constraints_propagated_sum}.
    Since $N\ge L$ and
    $L\ge (L_0+1)+(g-1)3(L_0+1)+1$, the product-span directness range of
    Lemma~\ref{lem:product_span_block_image_direct_sum} applies at length $N$:
    \begin{align}
        \spn\{(A^1_w,\ldots,A^g_w):|w|=N\}
         & =\prod_j \MN{D_j}
        \notag                  \\
         & \quad\Longrightarrow
        \ker\!\left(\bigoplus_{j=1}^g G_N(A_j)
        \longrightarrow(\C^d)^{\otimes N},
        \quad
        (\phi_j)_j\longmapsto\sum_j\phi_j\right)=0.
        \notag
    \end{align}
\end{proof}

\begin{lemma}[Periodic parent-space vectors decompose in the local block spaces]
    \label{lem:bnt_block_diagonal_open_boundary_decomposition}
    \lean{MPSTensor.exists_unique_sum_groundSpace_of_chainGroundSpace_toTensorFromBlocks_of_bnt_unital_c1}
    \leanok
    \uses{lem:bnt_block_diagonal_periodic_constraints_propagated_sum}
    With the hypotheses and notation of
    Lemma~\ref{lem:bnt_block_diagonal_periodic_constraints_propagated_sum},
    every vector
    $\psi\in\mathcal G_{N,L}(\bigoplus_{j=1}^g\mu_jA_j)$ has a unique family
    $(\psi_j)_{j=1}^g$ such that
    \begin{align}
        \psi   & =\sum_{j=1}^g\psi_j,
        \notag                        \\
        \psi_j & \in G_N(A_j).
        \notag
    \end{align}
\end{lemma}

\begin{proof}
    \leanok
    \uses{lem:bnt_block_diagonal_periodic_constraints_internal_sum}
    Lemma~\ref{lem:bnt_block_diagonal_periodic_constraints_internal_sum} gives
    \begin{align}
        \mathcal G_{N,L}\!\left(\bigoplus_{j=1}^g\mu_jA_j\right)
         & \subseteq\bigvee_{j=1}^gG_N(A_j),
        \notag
    \end{align}
    and the spaces $G_N(A_j)$ are independent: for every $j$ and every
    $x_j\in G_N(A_j)$,
    $\sum_jx_j=0 \quad\Longrightarrow\quad x_j=0$.
    Thus every $\psi\in\mathcal G_{N,L}(\bigoplus_j\mu_jA_j)$ has a unique
    decomposition $\psi=\sum_j\psi_j$ with $\psi_j\in G_N(A_j)$.
\end{proof}

\begin{lemma}[Block-diagonal boundary conditions for periodic parent-space vectors]
    \label{lem:bnt_block_diagonal_open_boundary_matrices}
    \lean{MPSTensor.exists_blockDiagonal_boundary_of_chainGroundSpace_toTensorFromBlocks_of_bnt_unital_c1}
    \leanok
    \uses{lem:bnt_block_diagonal_periodic_constraints_propagated_sum, lem:block_diagonal_boundary_of_open_boundary_sum}
    With the hypotheses and notation of
    Lemma~\ref{lem:bnt_block_diagonal_periodic_constraints_propagated_sum},
    every vector
    $\psi\in\mathcal G_{N,L}(\bigoplus_{j=1}^g\mu_jA_j)$ admits block-diagonal
    boundary conditions $X_j$ such that
    \begin{align}
        \psi
         & =\Gamma_N^B\!\left(\bigoplus_{j=1}^gX_j\right),
        \notag                                             \\
        \Gamma_N^{A_j}(\mu_j^NX_j)
         & \in G_N(A_j)
        \notag
    \end{align}
    for every $j$.
\end{lemma}

\begin{proof}
    \leanok
    \uses{lem:bnt_block_diagonal_periodic_constraints_propagated_sum, lem:block_diagonal_boundary_of_open_boundary_sum}
    Lemma~\ref{lem:bnt_block_diagonal_periodic_constraints_propagated_sum}
    places $\psi$ in $\bigvee_{j=1}^gG_N(A_j)$. Applying
    Lemma~\ref{lem:block_diagonal_boundary_of_open_boundary_sum} gives
    boundary matrices $X_j$ for which, for every $j$,
    $\Gamma_N^{A_j}(\mu_j^NX_j)\in G_N(A_j)$, and
    \begin{align}
        \psi
         & =\Gamma_N^B\!\left(\bigoplus_{j=1}^gX_j\right).
        \notag
    \end{align}
\end{proof}

\begin{lemma}[Block-diagonal boundary conditions at the source length]
    \label{lem:bnt_block_diagonal_open_boundary_matrices_source_length}
    \lean{MPSTensor.chainGroundSpace_toTensorFromBlocks_le_iSup_groundSpace_pgvwc07_of_dualFixedPoint}
    \lean{MPSTensor.chainGroundSpace_toTensorFromBlocks_le_iSup_and_iSupIndep_pgvwc07_of_dualFixedPoint}
    \lean{MPSTensor.exists_blockDiagonal_boundary_of_chainGroundSpace_toTensorFromBlocks_pgvwc07_of_dualFixedPoint}
    \leanok
    \uses{thm:pgvwc_direct_sum_intersection_source_length, thm:pgvwc_block_space_independence_source_length}
    Let $g\ge2$. Let $A_j$ be irreducible blocks, pairwise inequivalent up to
    gauge and phase, and injective at the common
    length $L_0>0$. Assume the canonical normalization of
    \cite{PerezGarcia2007Matrix}:
    \begin{align}
        \sum_aA^j_aA^{j\,\dagger}_a & =\Id,
                                    & \sum_aA^{j\,\dagger}_a\Lambda_jA^j_a & =\Lambda_j,
                                    & \Lambda_j                            & >0.
        \notag
    \end{align}
    Let every weight $\mu_j$ be nonzero. If
    $0<L\le N$ and $L\ge3(g-1)(L_0+1)+1$, then
    \begin{align}
        \mathcal G_{N,L}\!\left(\bigoplus_j\mu_jA_j\right)
         & \subseteq\bigoplus_jG_N(A_j).
        \notag
    \end{align}
    In particular, every vector $\psi$ in the left-hand side admits matrices
    $X_j$ such that
    \begin{align}
        \psi
         & =\Gamma_N^{\oplus_j\mu_jA_j}
        \left(\bigoplus_jX_j\right),
        \notag                          \\
        \Gamma_N^{A_j}(\mu_j^NX_j)
         & \in G_N(A_j).
        \notag
    \end{align}
    This is the open-boundary decomposition used in
    \cite[Theorem~12, proof lines~1424--1456]{PerezGarcia2007Matrix}.
\end{lemma}

\begin{proof}
    \leanok
    \uses{thm:pgvwc_direct_sum_intersection_source_length, thm:pgvwc_block_space_independence_source_length, lem:block_diagonal_boundary_of_open_boundary_sum}
    Put $S_M=\sum_jG_M(A_j)$. For every $M$ with $L\le M<N$,
    Theorem~\ref{thm:pgvwc_direct_sum_intersection_source_length} gives
    \begin{align}
        \left(\bigcap_b\Res_{-,b}^{-1}S_M\right)
        \cap
        \left(\bigcap_a\Res_{a,-}^{-1}S_M\right)
         & =S_{M+1}.
        \notag
    \end{align}
    Induction through the cyclic local constraints therefore gives
    \begin{align}
        \mathcal G_{N,L}\!\left(\bigoplus_j\mu_jA_j\right)
         & \subseteq S_N.
        \notag
    \end{align}
    Since $N\ge L\ge3(g-1)(L_0+1)+1$,
    Theorem~\ref{thm:pgvwc_block_space_independence_source_length} shows that
    the spaces $G_N(A_j)$ form an internal sum. Lemma~\ref{lem:block_diagonal_boundary_of_open_boundary_sum}
    then supplies matrices $X_j$ such that
    $\Gamma_N^{A_j}(\mu_j^NX_j)\in G_N(A_j)$ for every $j$ and
    \begin{align}
        \psi
         & =\Gamma_N^{\oplus_j\mu_jA_j}\!\left(\bigoplus_jX_j\right).
        \notag
    \end{align}
\end{proof}

\begin{theorem}[Doubly normalized block-diagonal boundary conditions at the source length]
    \label{thm:bnt_block_diagonal_open_boundary_matrices_source_length_doubly_normalized}
    \lean{MPSTensor.chainGroundSpace_toTensorFromBlocks_le_iSup_groundSpace_of_ge_of_bnt_directSum_unital_c1_pgvwc07}
    \lean{MPSTensor.exists_blockDiagonal_boundary_of_chainGroundSpace_toTensorFromBlocks_of_bnt_unital_c1_pgvwc07}
    \leanok
    Let $g\ge2$, assume the physical alphabet is nonempty, and let $A_j$
    be irreducible blocks with positive bond dimensions, pairwise inequivalent
    up to gauge and phase, and injective at the common length $L_0>0$. Assume
    normalized self-overlap and
    \begin{align}
        \sum_a(A^j_a)^\dagger A^j_a & =\Id,
        \notag                              \\
        \sum_aA^j_a(A^j_a)^\dagger  & =\Id.
        \notag
    \end{align}
    Let every weight $\mu_j$ be nonzero. If $0<L\le N$ and
    $L\ge3(g-1)(L_0+1)+1$, then
    \begin{align}
        \mathcal G_{N,L}\!\left(\bigoplus_j\mu_jA_j\right)
         & \subseteq\bigoplus_jG_N(A_j).
        \notag
    \end{align}
    In particular, every vector $\psi$ in the left-hand side admits matrices
    $X_j$ such that
    \begin{align}
        \psi
         & =\Gamma_N^{\oplus_j\mu_jA_j}\!\left(\bigoplus_jX_j\right),
        \notag                                                        \\
        \Gamma_N^{A_j}(\mu_j^NX_j)
         & \in G_N(A_j)
        \notag
    \end{align}
    for every $j$.
\end{theorem}

\begin{proof}\leanok
    \uses{thm:pgvwc_direct_sum_intersection_doubly_normalized, thm:pgvwc_block_space_independence_doubly_normalized, lem:block_diagonal_boundary_of_open_boundary_sum}
    Put $S_M=\sum_jG_M(A_j)$. For every $M$ with $L\le M<N$,
    Theorem~\ref{thm:pgvwc_direct_sum_intersection_doubly_normalized} gives
    \begin{align}
        \left(\bigcap_b\Res_{-,b}^{-1}S_M\right)
        \cap
        \left(\bigcap_a\Res_{a,-}^{-1}S_M\right)
         & =S_{M+1}.
        \notag
    \end{align}
    Induction through the cyclic local constraints therefore yields
    \begin{align}
        \mathcal G_{N,L}\!\left(\bigoplus_j\mu_jA_j\right)
         & \subseteq S_N.
        \notag
    \end{align}
    At the stated length, the block-space independence theorem makes this sum
    internal. Lemma~\ref{lem:block_diagonal_boundary_of_open_boundary_sum}
    then supplies the matrices $X_j$.
\end{proof}

The theorem gives $\Gamma_N^{A_j}(\mu_j^NX_j)\in G_N(A_j)$. This is not enough for
periodic-boundary membership when a cyclic window crosses the chosen cut. The
boundary-cut comparison of
Lemma~\ref{lem:block_diagonal_boundary_component_pgvwc_comparison_injective}
concerns the separate assertion
$\Gamma_N^{A_j}(\mu_j^NX_j)\in\mathcal G_{N,L}(A_j)$.

\begin{theorem}[Boundary compatibility in the trace-preserving normalization]
    \label{thm:ph_tp_boundary_compatibility}
    \lean{MPSTensor.boundary_matrix_eq_of_compatibility_of_tracePreserving}
    \leanok
    Suppose that $\sum_b (A^b)^\dagger A^b=\Id$ and that
    $A^b C_a=D_b A^a$ for every pair of physical letters $a,b$.
    Then $C_a=EA^a$, where $E=\sum_b(A^b)^\dagger D_b$.
    This is the trace-preserving form of the boundary calculation in
    \cite[Theorem~12]{PerezGarcia2007Matrix}.
\end{theorem}
\begin{proof}
    \leanok
    Multiplication on the left by $(A^b)^\dagger$ and summation over $b$
    gives $C_a=\sum_b(A^b)^\dagger D_b A^a=EA^a$.
\end{proof}

\begin{theorem}[Block intersections in the trace-preserving normalization]
    \label{thm:ph_tp_block_intersection}
    \lean{MPSTensor.mem_iSup_groundSpace_of_iSup_restrictions_of_tracePreserving,
        MPSTensor.iSup_groundSpace_eq_restriction_intersection_of_tracePreserving}
    \leanok
    Let $A^0,\ldots,A^{r-1}$ satisfy
    $\sum_a(A^{j,a})^\dagger A^{j,a}=\Id$ for each $j$.
    Suppose that the tuples $(A^{0,w},\ldots,A^{r-1,w})$, with $|w|=n$,
    span $\prod_j\MN{D_j}$. Writing $S_m=\sum_jG_m(A^j)$, one has
    \begin{align}
        (\C^d\otimes S_{n+1})\cap(S_{n+1}\otimes\C^d)=S_{n+2}.
        \notag
    \end{align}
\end{theorem}
\begin{proof}
    \leanok
    \uses{lem:pgvwc07_trace_decompositions_from_block_restrictions,
        lem:blockwise_boundary_compatibility_from_traces,
        thm:ph_tp_boundary_compatibility,
        lem:left_boundary_summands_mem_block_ground_spaces,
        lem:ground_space_in_left_ground, lem:ground_space_in_right_ground}
    The two restrictions provide boundary matrices $C_a^j,D_b^j$
    (Lemma~\ref{lem:pgvwc07_trace_decompositions_from_block_restrictions}).
    The simultaneous word span gives $A^{j,b}C_a^j=D_b^j A^{j,a}$
    (Lemma~\ref{lem:blockwise_boundary_compatibility_from_traces}).
    Set $E_j=\sum_b(A^{j,b})^\dagger D_b^j$. Then $C_a^j=E_jA^{j,a}$ by
    Theorem~\ref{thm:ph_tp_boundary_compatibility}, so each boundary summand
    lies in $G_{n+2}(A^j)$
    (Lemma~\ref{lem:left_boundary_summands_mem_block_ground_spaces}).
    The reverse inclusion follows by restricting a boundary-condition vector
    at its first or last physical letter
    (Lemmas~\ref{lem:ground_space_in_left_ground}
    and~\ref{lem:ground_space_in_right_ground}):
    \begin{align}
        G_{n+2}(A^j)\subseteq\C^d\otimes G_{n+1}(A^j),
        \qquad
        G_{n+2}(A^j)\subseteq G_{n+1}(A^j)\otimes\C^d.
        \notag
    \end{align}
\end{proof}

\begin{theorem}[Open-chain ground space of trace-preserving blocks]
    \label{thm:ph_tp_block_open_ground_space}
    \lean{MPSTensor.contiguous_mem_groundSpace_toTensorFromBlocks,
        MPSTensor.ker_openParentHamiltonianES_toTensorFromBlocks_eq_groundSpaceES}
    \leanok
    Let $d>0$, and let $B=\bigoplus_j\mu_j A^j$, where every $\mu_j$ is nonzero
    and $\sum_a(A^{j,a})^\dagger A^{j,a}=\Id$ for every $j$.
    Suppose that the simultaneous word tuples span the full product matrix
    algebra at every length $n\ge L_0$. For $L_0+1\le L\le N$, the
    open-chain parent Hamiltonian with interaction range $L$ satisfies
    \begin{align}
        \ker H^{\mathrm{open}}_{B,L,N}=G_N(B)=\sum_jG_N(A^j).
        \notag
    \end{align}
\end{theorem}
\begin{proof}
    \leanok
    \uses{thm:ph_tp_block_intersection,
        lem:open_chain_iteration_subspace_sequence,
        thm:block_sum_local_ground_space, def:open_parent_hamiltonian_es,
        thm:open_parent_hamiltonian_zero_local_terms,
        lem:local_term_es_kernel_restrictions, lem:mem_ground_space_es,
        thm:open_ground_space_le_hamiltonian_kernel}
    By Definition~\ref{def:open_parent_hamiltonian_es}, the Hamiltonian is a
    sum of orthogonal projectors,
    \begin{align}
        H^{\mathrm{open}}_{B,L,N}=\sum_{i\in I_{L,N}}h_{i,\mathrm{ES}}(B,L),
        \qquad h_{i,\mathrm{ES}}(B,L)\ge0.
        \notag
    \end{align}
    Hence, by Theorem~\ref{thm:open_parent_hamiltonian_zero_local_terms},
    \begin{align}
        \ker H^{\mathrm{open}}_{B,L,N}
        =\bigcap_{i\in I_{L,N}}\ker h_{i,\mathrm{ES}}(B,L),
        \notag
    \end{align}
    so, by Lemma~\ref{lem:local_term_es_kernel_restrictions}, every
    contiguous length-$L$ restriction of a zero-energy vector lies in
    $G_L(B)$ under the identification of
    Lemma~\ref{lem:mem_ground_space_es}. The block intersection
    identity of Theorem~\ref{thm:ph_tp_block_intersection} extends these
    constraints one site at a time
    (Lemma~\ref{lem:open_chain_iteration_subspace_sequence}), giving
    membership in $G_N(B)$. Conversely,
    $G_N(B)\subseteq\ker H^{\mathrm{open}}_{B,L,N}$ by
    Theorem~\ref{thm:open_ground_space_le_hamiltonian_kernel}. Nonzero
    block weights identify $G_N(B)$ with the sum of the block spaces
    (Theorem~\ref{thm:block_sum_local_ground_space}).
\end{proof}

\begin{theorem}[Three-block pairwise separation from block injectivity]
    \label{thm:ph_three_block_pairwise_separation}
    \lean{MPSTensor.hasPairBlockSeparatingWords_threeBlock_of_blocksNotGaugePhaseEquiv_c1}
    \leanok
    \uses{def:pairwise_block_separating_words, def:pair_trace_separation_words,
        def:l_blk_injective, def:blocks_not_gauge_phase_equiv}
    Let $A^0,\ldots,A^{r-1}$ be tensors with positive bond dimensions such
    that each $A^j$ is irreducible, satisfies
    $\sum_a(A^{j,a})^\dagger A^{j,a}=\Id$, and has self-overlap tending to
    $1$ with the chain length. Suppose that no two distinct blocks are
    related by a virtual gauge and a nonzero complex scalar. Let $L_0>0$ and
    suppose that every $A^j$ is block injective at the lengths $L_0$,
    $L_0+1$ and $S=3(L_0+1)$. Then every ordered pair of distinct blocks is
    trace separated at length $S$, and for all $k\ne j$ there are
    coefficients $c_\omega$ with
    \begin{align}
        \sum_{|\omega|=S}c_\omega A_k^\omega=I,
        \qquad
        \sum_{|\omega|=S}c_\omega A_j^\omega=0.
        \notag
    \end{align}
\end{theorem}
\begin{proof}
    \leanok
    \uses{lem:bnt_direct_sum_injective_prefix_word_span,
        lem:direct_sum_two_block_dimension_step,
        lem:pair_trace_separation_to_pairwise_block_separation}
    Fix $k\ne j$. If $D_k\ne D_j$, the strict-size branch of
    Lemma~\ref{lem:direct_sum_two_block_dimension_step} at the comparison
    length $L_0+1$ gives
    \begin{align}
        \spn\{(A_k^\omega,A_j^\omega):|\omega|=S\}=\MN{D_k}\times\MN{D_j}.
        \notag
    \end{align}
    If $D_k=D_j$, the blocks are not related by a virtual gauge and a nonzero
    scalar, so their periodic matrix product vectors are not proportional at
    arbitrarily long chain lengths, and the same lemma gives
    $\mathcal G_S^{A_k}\cap\mathcal G_S^{A_j}=\{0\}$ and the same pair span.
    In both cases the pair is trace separated at length $S$, which is
    Lemma~\ref{lem:bnt_direct_sum_injective_prefix_word_span}, and
    Lemma~\ref{lem:pair_trace_separation_to_pairwise_block_separation}
    converts the pair span into the displayed equations.
\end{proof}

\begin{theorem}[Eventual word span from block selectors]
    \label{thm:ph_eventual_word_span_from_selectors}
    \lean{MPSTensor.eventually_wordTupleSpanTop_of_blockSelectorWords_of_isNBlkInjective}
    \leanok
    \uses{def:l_blk_injective, def:blockwise_product_word_span}
    Let $A^0,\ldots,A^{r-1}$ be tensors, let $p>0$, and suppose that every
    $A^j$ is block injective at length $p$. Suppose that for every $k$ there
    are coefficients $c_\omega^{(k)}$ with
    \begin{align}
        \sum_{|\omega|=s}c_\omega^{(k)}A_j^\omega=
        \begin{cases}
            I, & j=k,    \\
            0, & j\ne k.
        \end{cases}
        \notag
    \end{align}
    Then there is $L_0$ such that, for every $L\ge L_0$, the tuples
    $(A^{0,w},\ldots,A^{r-1,w})$, with $|w|=L$, span $\prod_j\MN{D_j}$.
\end{theorem}
\begin{proof}
    \leanok
    \uses{thm:wordspan_blk_inj_succ,
        lem:product_word_span_from_bnt_block_separation}
    Take $L_0=s+p$ and let $L\ge L_0$. Since $L-s\ge p$, every $A^j$ is
    block injective at length $L-s$ by Theorem~\ref{thm:wordspan_blk_inj_succ}. Lemma
    \ref{lem:product_word_span_from_bnt_block_separation} with the selectors
    of length $s$ gives the full product span at length $(L-s)+s=L$.
\end{proof}

\begin{theorem}[Eventual simultaneous word span of primitive blocks]
    \label{thm:ph_primitive_blocks_word_span}
    \lean{MPSTensor.exists_eventually_wordTupleSpanTop_of_isPrimitiveMPS}
    \leanok
    Let $A^0,\ldots,A^{r-1}$ be trace-preserving primitive tensors with
    positive bond dimensions and positive definite invariant matrices.
    Suppose that no two distinct blocks are related by conjugation and a
    nonzero complex scalar. Then there is $L_0$ such that, for every $L\ge L_0$,
    the tuples $(A^{0,w},\ldots,A^{r-1,w})$, with $|w|=L$, span
    $\prod_j\MN{D_j}$.
\end{theorem}
\begin{proof}
    \leanok
    \uses{thm:normal_from_primitive_posdef,
        thm:common_exact_block_injectivity_normal_left_canonical,
        thm:wordspan_blk_inj_succ,
        thm:ph_three_block_pairwise_separation,
        lem:block_separating_equations_from_pairwise_separation,
        thm:ph_eventual_word_span_from_selectors}
    Primitivity gives normality
    (Theorem~\ref{thm:normal_from_primitive_posdef}), hence a common positive
    length $p$ at which every block is injective
    (Theorem~\ref{thm:common_exact_block_injectivity_normal_left_canonical}).
    Injectivity persists at every length at least $p$
    (Theorem~\ref{thm:wordspan_blk_inj_succ}), in particular at $p+1$ and
    $3(p+1)$. Theorem~\ref{thm:ph_three_block_pairwise_separation} with
    $L_0=p$ gives, for every ordered pair $k\ne j$, a word polynomial of
    length $3(p+1)$ equal to $I$ on block $k$ and to $0$ on block $j$.
    Lemma~\ref{lem:block_separating_equations_from_pairwise_separation}
    multiplies these polynomials into selectors of the common length
    $s=(r-1)\cdot3(p+1)$:
    \begin{align}
        \sum_{|\omega|=s}c_\omega^{(k)}A_j^\omega=
        \begin{cases}
            I, & j=k,    \\
            0, & j\ne k.
        \end{cases}
        \notag
    \end{align}
    Theorem~\ref{thm:ph_eventual_word_span_from_selectors} then gives the full
    product span at every length $L\ge s+p$.
\end{proof}

\begin{theorem}[Uniform open-chain kernel threshold for primitive blocks]
    \label{thm:ph_primitive_blocks_open_kernel}
    \lean{MPSTensor.exists_ker_openParentHamiltonianES_toTensorFromBlocks_eq_groundSpaceES_of_isPrimitiveMPS}
    \leanok
    Let $d>0$ and $B=\bigoplus_j\mu_jA^j$, with nonzero coefficients $\mu_j$.
    Suppose that the blocks have positive bond dimensions, are trace-preserving
    and primitive with positive definite invariant matrices, and are pairwise
    inequivalent under conjugation and a nonzero scalar.
    There is $L_0>0$ such that for every $L_0\le L\le N$,
    \begin{align}
        \ker H^{\mathrm{open}}_{B,L,N}=G_N(B).
        \notag
    \end{align}
\end{theorem}
\begin{proof}
    \leanok
    \uses{thm:ph_primitive_blocks_word_span,
        thm:ph_tp_block_open_ground_space}
    Choose $L_0$ one larger than the eventual simultaneous word-span
    threshold of Theorem~\ref{thm:ph_primitive_blocks_word_span} and apply
    Theorem~\ref{thm:ph_tp_block_open_ground_space}.
\end{proof}

\begin{theorem}[Primitive representatives of inequivalent normal blocks]
    \label{thm:normal_blocks_primitive_representatives}
    \lean{MPSTensor.exists_isPrimitiveMPS_family_of_isNormal}
    \leanok
    \uses{def:normal, def:primitive_mps, def:ground_space_parent,
        def:gauge_phase_equiv}
    Let $(A^i)_{i\in I}$ be normal tensors with positive bond dimensions.
    Suppose that no two distinct blocks are related by a virtual gauge
    and a nonzero complex scalar. Then there are tensors $(B^i)_{i\in I}$
    of the same respective bond dimensions and positive definite matrices
    $\rho_i$ such that each $B^i$ is primitive and trace preserving,
    $\rho_i$ is its invariant matrix, and $G_L(A^i)=G_L(B^i)$ for every
    $i$ and every length $L$. The tensors $B^i$ remain pairwise inequivalent
    under a virtual gauge and a nonzero scalar.
\end{theorem}
\begin{proof}
    \leanok
    \uses{thm:normalized_primitive_gauge,
        lem:gauge_phase_equiv_transport_left_right}
    Normalize each block independently by
    Theorem~\ref{thm:normalized_primitive_gauge}. The resulting nonzero
    rescaling and virtual gauge leave its local ground spaces unchanged.
    Thus $B^{i,a}=\zeta_iX_iA^{i,a}X_i^{-1}$ for every physical index $a$,
    with $\zeta_i\ne0$ and $X_i$ invertible. If $D_i=D_j$ and
    $B^{i,a}=\zeta YB^{j,a}Y^{-1}$ for every $a$, with $\zeta\ne0$ and $Y$
    invertible, then composing the gauges gives
    \begin{align}
        A^{i,a}=\frac{\zeta\zeta_j}{\zeta_i}\,
        \bigl(X_i^{-1}YX_j\bigr)A^{j,a}\bigl(X_i^{-1}YX_j\bigr)^{-1}
        \notag
    \end{align}
    for every $a$ (Lemma~\ref{lem:gauge_phase_equiv_transport_left_right}),
    contrary to the hypothesis.
\end{proof}

\begin{theorem}[Blockwise replacement preserves parent Hamiltonians]
    \label{thm:ph_block_replacement_parent_equality}
    \lean{MPSTensor.groundSpace_toTensorFromBlocks_eq_of_block_groundSpace_eq,
        MPSTensor.parentHamiltonianES_toTensorFromBlocks_eq_of_block_groundSpace_eq}
    \leanok
    Let $A^j$ and $B^j$ have the same respective bond dimensions, and let
    all coefficients $\mu_j,\nu_j$ be nonzero. If $G_L(A^j)=G_L(B^j)$ for
    every $j$, then the assembled tensors $A=\bigoplus_j\mu_jA^j$ and
    $B=\bigoplus_j\nu_jB^j$ satisfy $G_L(A)=G_L(B)$.
    Their canonical periodic parent Hamiltonians of range $L$ coincide at
    every volume: $H_{A,L,N}=H_{B,L,N}$.
\end{theorem}
\begin{proof}
    \leanok
    \uses{thm:block_sum_local_ground_space}
    The assembled local ground spaces are the sums of the respective block
    spaces. Equality of these spaces gives equality of their orthogonal
    complement projections and hence of every translated parent interaction.
\end{proof}

\begin{theorem}[Primitive block representatives with unchanged Hamiltonians]
    \label{thm:ph_primitive_representatives_parent_equality}
    \lean{MPSTensor.exists_isPrimitiveMPS_family_parentHamiltonianES_eq_of_isNormal}
    \leanok
    Let $A=\bigoplus_j\mu_jA^j$ have nonzero coefficients and normal blocks
    of positive bond dimensions. Suppose that no two distinct blocks are
    related by conjugation and a nonzero scalar.
    There are normalized primitive blocks $B^j$, with positive definite
    invariant matrices, which remain pairwise inequivalent and satisfy
    $H_{A,L,N}=H_{B,L,N}$ for every $L,N$, where
    $B=\bigoplus_j\mu_jB^j$.
\end{theorem}
\begin{proof}
    \leanok
    \uses{thm:normal_blocks_primitive_representatives,
        thm:ph_block_replacement_parent_equality}
    Normalize each block while retaining its local ground spaces and the
    inequivalence of distinct blocks. The assembled parent Hamiltonians
    therefore coincide at each interaction range and volume.
\end{proof}

\begin{lemma}[Primitive normalization at a fixed C1 length]
    \label{lem:normal-block-primitive-c1-normalization}
    \lean{MPSTensor.exists_isPrimitiveMPS_family_parentHamiltonianES_eq_of_common_isNBlkInjective}
    \leanok
    Suppose that a finite family of nonzero-dimensional MPS tensors is
    injective at a common positive word length $L_0$, and that its members
    are pairwise inequivalent under virtual conjugation and nonzero scalar
    multiplication. Then each tensor admits a normalized primitive
    representative which remains injective at length $L_0$.
    The representatives remain pairwise inequivalent. For any nonzero
    weights, the original and normalized block sums have exactly the same
    periodic parent Hamiltonians at every interaction range and chain length.
\end{lemma}
\begin{proof}
    \leanok
    \uses{thm:normalized_primitive_gauge,
        thm:ph_block_replacement_parent_equality}
    By Theorem~\ref{thm:normalized_primitive_gauge}, normalize each tensor by
    a nonzero scalar and an invertible virtual conjugation. These operations
    preserve injectivity at each fixed word length, preserve all local
    ground spaces, and preserve inequivalence. By
    Theorem~\ref{thm:ph_block_replacement_parent_equality}, equal block
    ground spaces give equal ground spaces of the weighted block sums, hence
    every local parent projection, and therefore every periodic parent
    Hamiltonian, is unchanged.
\end{proof}

\begin{theorem}[Propagation of simultaneous spans in trace-preserving form]
    \label{thm:ph_tp_word_span_propagation}
    \lean{MPSTensor.wordTupleSpanTop_succ_of_tracePreserving,
        MPSTensor.wordTupleSpanTop_of_ge_of_tracePreserving}
    \leanok
    Suppose that $\sum_a(A^{j,a})^\dagger A^{j,a}=\Id$ for every block $j$.
    If the simultaneous length-$L$ word tuples span the product matrix
    algebra, the same holds at every length $n\ge L$.
\end{theorem}
\begin{proof}
    \leanok
    For an arbitrary tuple $(M_j)_j$, represent
    $(M_j(A^{j,a})^\dagger)_j$ as a linear combination of length-$L$ word
    tuples. Append the letter $a$ and sum over $a$. The resulting tuple is
    $(M_j)_j$ by trace-preserving normalization. This proves the one-site
    extension, and induction gives every larger length.
\end{proof}

\begin{theorem}[The printed open-chain threshold for primitive blocks]
    \label{thm:ph_primitive_blocks_sharp_open_kernel}
    \lean{MPSTensor.ker_openParentHamiltonianES_toTensorFromBlocks_eq_groundSpaceES_of_threeBlock_bound}
    \leanok
    Let $d>0$, $r\ge2$, and $B=\bigoplus_{j=1}^r\mu_jA^j$, with all
    coefficients nonzero. Suppose that the blocks have positive bond
    dimensions, are trace-preserving and primitive with positive definite
    invariant matrices, and are pairwise inequivalent under conjugation
    and multiplication by a nonzero scalar.
    Suppose that each block is injective at the supplied common length
    $L_0>0$. Then for
    \begin{align}
        3(r-1)(L_0+1)+1\le R\le N,
        \notag
    \end{align}
    the open-chain parent Hamiltonian satisfies
    $\ker H^{\mathrm{open}}_{B,R,N}=G_N(B)$.
    This is the interaction threshold of
    \cite[Theorem~12]{PerezGarcia2007Matrix} in trace-preserving coordinates.
\end{theorem}
\begin{proof}
    \leanok
    \uses{thm:ph_tp_word_span_propagation,
        thm:ph_tp_block_open_ground_space}
    The simultaneous-span estimate of
    \cite[Lemma~10]{PerezGarcia2007Matrix} gives the full product matrix
    algebra at length $3(r-1)(L_0+1)$. Trace-preserving propagation keeps
    the span full at every larger length. The block intersection theorem
    therefore gives the kernel identity at the stated interaction range.
\end{proof}

\begin{theorem}[A periodic gap for normal block sums]
    \label{thm:normal_block_sum_parent_gap}
    \lean{MPSTensor.exists_parentHamiltonianES_toTensorFromBlocks_gap_of_isNormal, MPSTensor.exists_parentHamiltonianES_toTensorFromBlocks_gap_threshold_of_isNormal}
    \leanok
    Let $A^0,\ldots,A^{r-1}$ be normal MPS tensors with positive physical and
    bond dimensions, pairwise inequivalent under virtual conjugation and
    nonzero scalar multiplication. Let $B=\bigoplus_j\mu_j A^j$, where
    every $\mu_j$ is nonzero.
    There is a positive integer $R_0$ such that, for every $R\geq R_0$,
    there are constants $\gamma_R>0$ and $N_R\in\mathbb N$ for which
    the canonical periodic parent Hamiltonian $H_{R,N}$ satisfies
    \begin{align}
        \gamma_R\lVert v\rVert\leq\lVert H_{R,N}v\rVert
    \end{align}
    for every $N\geq N_R$ and every $v\perp\ker H_{R,N}$.
    In particular, at least one positive interaction range has a uniform
    eventual periodic gap.
\end{theorem}
\begin{proof}\leanok
    \uses{thm:ph_primitive_representatives_parent_equality, thm:primitive_block_sum_parent_gap_threshold}
    Replace each block by a normalized primitive representative. This preserves
    inequivalence and leaves every parent Hamiltonian unchanged. The uniform
    gap theorem for primitive block sums therefore applies with the same
    interaction ranges and the same constants.
\end{proof}

\begin{theorem}[The gap at the prescribed interaction range]
    \label{thm:normal-block-parent-gap-three-block-bound}
    \lean{MPSTensor.exists_parentHamiltonianES_toTensorFromBlocks_uniform_gap_of_common_isNBlkInjective}
    \leanok
    Let $A^0,\ldots,A^{r-1}$, $r\geq2$, be MPS tensors with positive
    physical and bond dimensions, injective at a common positive word length
    $L_0$, and pairwise inequivalent under virtual conjugation and nonzero
    scalar multiplication. Let $B=\bigoplus_j\mu_j A^j$, with every
    $\mu_j\ne0$. For every prescribed range
    \begin{align}
        R\geq3(r-1)(L_0+1)+1,
    \end{align}
    there is a constant $\gamma_R>0$ such that the canonical periodic parent
    Hamiltonian $H_{R,N}$ of $B$ satisfies
    \begin{align}
        \gamma_R\lVert v\rVert\leq\lVert H_{R,N}v\rVert
    \end{align}
    for every chain length $N\in\mathbb N$ and every $v\perp\ker H_{R,N}$.
    This is the block-injective parent-Hamiltonian gap of
    \cite[Section~IV.C]{Cirac2021Matrix} at the interaction bound of
    \cite[Theorem~12]{PerezGarcia2007Matrix}.
\end{theorem}
\begin{proof}\leanok
    \uses{lem:normal-block-primitive-c1-normalization, thm:primitive-block-parent-uniform-gap-three-block-bound}
    Normalize the blocks to primitive representatives without changing the
    specified injectivity length $L_0$. Their periodic parent Hamiltonians
    agree with the original ones at every range and chain length. Apply the
    primitive-block gap theorem at the prescribed range.
\end{proof}

\begin{theorem}[Cyclic propagation of boundary intertwiners]
    \label{thm:cyclic-boundary-intertwining}
    \lean{MPSTensor.boundary_commutes_evalWord_of_cyclic_intertwining}
    \leanok
    Let $N\geq1$, let $s\in\mathbb Z/N\mathbb Z$, and suppose matrices $Y_i$
    satisfy $Y_iA^a=A^aY_{s+i}$ for every $i$ and physical label $a$.
    Then, for every word $w$ of length $N$ and every $i$,
    \begin{align}
        Y_iA^w=A^wY_i.
    \end{align}
\end{theorem}
\begin{proof}\leanok
    \uses{thm:cyclic_window_boundary_matrix_transport}
    Propagating the intertwining identity through the word gives
    $Y_iA^w=A^wY_{Ns+i}$. Since $Ns=0$ in the cyclic index group, the two
    boundary matrices coincide.
\end{proof}

\begin{theorem}[Cyclic closure of block boundary matrices]
    \label{thm:ph_block_cyclic_boundary_closure}
    \lean{MPSTensor.mem_bntMPSVectorSpan_of_cyclic_openBoundary}
    \leanok
    \uses{def:parent_cyclic_translate_state, def:ground_space_map}
    Let $N\geq2$. Suppose that the simultaneous words of length $N-1$
    in the tensors $A^j$ span the product of their matrix algebras.
    Let $\psi$ be a vector such that, for every cyclic cut
    $i\in\mathbb Z/N\mathbb Z$, there are boundary matrices $X_i^j$ with
    \begin{align}
        T_i\psi=\sum_j\Gamma_N^{A^j}(X_i^j).
        \notag
    \end{align}
    Then $\psi$ belongs to the span of the periodic MPS vectors of the
    blocks.
    This is the block-diagonal cyclic closure argument in
    \cite[Section~IV.C]{Cirac2021Matrix}.
\end{theorem}
\begin{proof}
    \leanok
    \uses{lem:block_boundary_intertwiner_from_global_cut,
        thm:cyclic-boundary-intertwining,
        thm:ground_space_map_mpv_of_long_word_commutes,
        lem:bnt_span_eq_sum_mpv_submodules}
    Choose the boundary matrices $X_i^j$ of the hypothesis at each cyclic
    cut $i$. Comparing adjacent cuts against all complementary words of length $N-1$
    gives $X_i^j A^{j,a}=A^{j,a}X_{i-1}^j$.
    Iterating this identity once around the ring yields
    \begin{align}
        X_i^j A^{j,a_1}\cdots A^{j,a_N}
        =A^{j,a_1}\cdots A^{j,a_N}X_i^j.
        \notag
    \end{align}
    Each block is injective at length $N-1$, hence also at length $N$.
    Its boundary matrix therefore commutes with the full matrix algebra
    and is scalar. The representation at cut zero is consequently a
    linear combination of the periodic block vectors, that is, an element of
    the sum of their lines
    (Lemma~\ref{lem:bnt_span_eq_sum_mpv_submodules}).
\end{proof}

\begin{theorem}[Periodic ground space at a simultaneous injectivity length]
    \label{thm:ph_block_periodic_closure_tp}
    \lean{MPSTensor.ker_parentHamiltonian_toTensorFromBlocks_eq_of_wordTupleSpanTop_of_tracePreserving}
    \leanok
    Let $B=\bigoplus_j\mu_j A^j$ have nonzero block coefficients and positive
    physical dimension. Suppose that each block is trace preserving and
    that the simultaneous words of length $S>0$ span the product matrix
    algebra. For every $S+1\leq L\leq N$, the periodic parent Hamiltonian
    satisfies
    \begin{align}
        \ker H_L^{(N)}(B)=\operatorname{span}\{V^{(N)}(A^j):j\}.
        \notag
    \end{align}
\end{theorem}
\begin{proof}
    \leanok
    \uses{thm:ph_block_cyclic_boundary_closure,thm:ph_tp_word_span_propagation,
        thm:ph_tp_block_open_ground_space, lem:bnt_span_le_parent_hamiltonian_kernel}
    Trace-preserving normalization propagates simultaneous injectivity to
    every length at least $S$, in particular to $N-1$. The open-chain
    intersection property expresses each periodic ground vector as a sum
    of block boundary vectors at every cut. Cyclic closure places this
    vector in the span of the periodic block vectors. Conversely, each
    periodic block vector satisfies every local parent constraint.
\end{proof}

\begin{theorem}[Independent rescaling of simultaneous word spans]
    \label{thm:block_word_span_independent_rescaling}
    \lean{MPSTensor.wordTupleSpanTop_of_family_smul}
    \leanok
    Suppose the length-$S$ word tuples of a finite family $A^0,\ldots,A^{r-1}$
    span the full product matrix algebra. If every $\zeta_j$ is nonzero,
    the word tuples of $\zeta_0A^0,\ldots,\zeta_{r-1}A^{r-1}$ also span that
    algebra at the same length $S$.
\end{theorem}
\begin{proof}\leanok
    Rescaling the $j$th block multiplies its length-$S$ words by
    $\zeta_j^S$. For any target tuple $(M_j)_j$, first express
    $(\zeta_j^{-S}M_j)_j$ in the original word span. The same coefficients
    express $(M_j)_j$ in the rescaled word span.
\end{proof}

\begin{theorem}[Primitive normalization at a simultaneous injectivity length]
    \label{thm:block_word_span_primitive_normalization}
    \lean{MPSTensor.exists_isPrimitiveMPS_family_of_wordTupleSpanTop, MPSTensor.exists_isPrimitiveMPS_family_parentHamiltonian_eq_of_wordTupleSpanTop}
    \leanok
    Let $S>0$, and suppose the length-$S$ word tuples of a finite family
    $A^0,\ldots,A^{r-1}$ with positive bond dimensions span the full product
    matrix algebra. There are normalized primitive representatives
    $B^0,\ldots,B^{r-1}$, with positive definite fixed points, whose word tuples
    still span the full product algebra at length $S$.
    For every $0\leq j<r$, every length $L$, and every chain length $N$,
    \begin{align}
        G_L(A^j) & =G_L(B^j),\notag                                      \\
        \operatorname{span}\{V^{(N)}(A^k):0\leq k<r\}
                 & =\operatorname{span}\{V^{(N)}(B^k):0\leq k<r\}.\notag
    \end{align}
    For any nonzero weights $\mu_j$, the periodic parent Hamiltonians
    of $\bigoplus_j\mu_j A^j$ and $\bigoplus_j\mu_j B^j$ are equal at every
    interaction range and chain length.
\end{theorem}
\begin{proof}\leanok
    \uses{thm:block_word_span_independent_rescaling, thm:word_tuple_span_gauge_invariance,
        thm:normalized_primitive_gauge, thm:block_sum_local_ground_space,
        lem:bnt_span_eq_sum_mpv_submodules}
    Projection of the simultaneous span onto each matrix summand proves
    injectivity of every block at length $S$. Each block is therefore normal
    and admits a primitive representative obtained by a rescaling by some
    $\zeta_j\neq0$ and an invertible virtual conjugation. Both operations
    preserve the simultaneous word span at the specified length. They
    preserve every component ground space, and the periodic component
    vectors satisfy
    \begin{align}
        V^{(N)}(B^j) & =\zeta_j^{\,N}\,V^{(N)}(A^j).\notag
    \end{align}
    Since $\zeta_j^{\,N}\neq0$, each block spans the same line for $A^j$ and
    $B^j$, and the two spans, being the sums of these lines
    (Lemma~\ref{lem:bnt_span_eq_sum_mpv_submodules}), agree. Taking sums of the
    component ground spaces proves equality of the assembled parent
    Hamiltonians.
\end{proof}

\begin{theorem}[Periodic ground space at simultaneous block injectivity]
    \label{thm:ph_block_ground_space_at_simultaneous_injectivity}
    \lean{MPSTensor.chainGroundSpace_toTensorFromBlocks_eq_of_wordTupleSpanTop, MPSTensor.ker_parentHamiltonian_toTensorFromBlocks_eq_of_wordTupleSpanTop, MPSTensor.ker_sum_eq_bntMPSVectorSpan_of_wordTupleSpanTop}
    \leanok
    Let $A^0,\ldots,A^{r-1}$ have a common positive physical dimension and
    positive bond dimensions. Suppose their length-$S$ word tuples span the
    full product matrix algebra, where $S>0$, and let
    $B=\bigoplus_j\mu_j A^j$ with all $\mu_j\ne0$.
    For every interaction range and chain length satisfying
    \begin{align}
        S+1\leq L\leq N,
        \notag
    \end{align}
    the periodic local constraints have ground space
    \begin{align}
        G_L^{(N)}(B)
        =\operatorname{span}\{V^{(N)}(A^j):0\leq j<r\}.
        \notag
    \end{align}
    This is the kernel of the canonical parent Hamiltonian. It is also the
    kernel of the sum of any positive terms whose cyclic local kernels agree
    with those of the canonical parent terms.
    The result includes the minimal ring $N=S+1$ and gives the
    block-injective closure statement of \cite[Section~IV.C]{Cirac2021Matrix}
    at the interaction range specified by the preceding intersection argument.
\end{theorem}
\begin{proof}\leanok
    \uses{thm:block_word_span_primitive_normalization, thm:ph_block_periodic_closure_tp,
        thm:parent_hamiltonian_kernel_eq_chain_ground_space}
    Normalize the blocks while retaining simultaneous injectivity at length
    $S$. The parent Hamiltonians and the span of the periodic component
    vectors are unchanged. Apply the trace-preserving block closure theorem.
    Finally, the kernel of a sum of positive operators is the intersection
    of their kernels, so replacing the local projectors by positive terms
    with the same local kernels preserves the ground space.
\end{proof}

\begin{theorem}[Local spaces of canonical representatives]
    \label{thm:ph_canonical_representative_local_spaces}
    \lean{MPSTensor.CPSVCanonicalFormData.groundSpace_eq_iSup_representatives}
    \leanok
    Let $A$ be in canonical form, and let $B^0,\ldots,B^{g-1}$ be its distinct
    normal representatives. Repeated copies and their nonzero coefficients
    do not change the positive-length local spaces: for every $L>0$,
    \begin{align}
        G_L(A)=\bigvee_{j=0}^{g-1}G_L(B^j).
        \notag
    \end{align}
    This is the local-space consequence of the canonical decomposition in
    \cite{Cirac2021Matrix}.
\end{theorem}
\begin{proof}
    \leanok
    \uses{lem:coisometric_reconstruction_local_space, thm:block_sum_local_ground_space,
        thm:local_ground_space_gauge_invariance, thm:local_ground_space_rescaling}
    The coisometric reconstruction preserves the local space on every
    positive-length interval. The local space of the resulting direct sum
    is the sum of its block spaces. Each repeated copy differs from its
    representative by a virtual similarity and a nonzero scalar, both of
    which preserve the local space. Taking the sum once over each distinct
    representative therefore gives the same space.
\end{proof}

\begin{theorem}[Canonical ground spaces at a simultaneous injectivity length]
    \label{thm:ph_canonical_block_closure}
    \lean{MPSTensor.CPSVCanonicalFormData.ker_parentHamiltonian_eq_of_wordTupleSpanTop,
        MPSTensor.CPSVCanonicalFormData.ker_sum_eq_of_wordTupleSpanTop,
        MPSTensor.IsCPSVCanonicalForm.exists_bnt_ker_sum_eq_of_wordTupleSpanTop}
    \leanok
    Let $A$ be a canonical MPS tensor with positive physical dimension,
    including arbitrary multiplicities of its distinct normal blocks
    $B^0,\ldots,B^{g-1}$. Suppose that their simultaneous words of length
    $S>0$ span the product matrix algebra. For every $S+1\leq L\leq N$,
    every periodic parent Hamiltonian with positive interactions and the
    prescribed length-$L$ local kernels satisfies
    \begin{align}
        \ker H_L^{(N)}(A)=\operatorname{span}\{V^{(N)}(B^j):0\leq j<g\}.
        \notag
    \end{align}
    The conclusion includes the minimal ring $N=L=S+1$ and the canonical
    orthogonal-projection interaction. This is the block-injective closure
    conclusion of \cite[Section~IV.C]{Cirac2021Matrix}.
\end{theorem}
\begin{proof}
    \leanok
    \uses{thm:ph_canonical_representative_local_spaces,
        thm:ph_block_ground_space_at_simultaneous_injectivity,
        thm:parent_hamiltonian_kernel_eq_chain_ground_space,
        thm:block_sum_local_ground_space}
    Replace the original canonical tensor by the direct sum containing
    one copy of each normal representative. Their local spaces, hence
    their periodic local constraints, coincide. Apply the simultaneous
    block-injectivity closure theorem at the same supplied length $S$.
    Finally, the kernel of a finite sum of positive operators is the
    intersection of their kernels. Thus replacing each orthogonal
    projection by any positive interaction with the same kernel leaves
    the periodic ground space unchanged.
\end{proof}
